% ============================================================================
%  Problems: เรขาคณิตวิเคราะห์และเวกเตอร์ (Analytic Geometry & Vectors)
%  Ordered from easy (basic) → intermediate.
%  Shared by exercise.tex and answer-key.tex
% ============================================================================

% --- Part 1: พื้นฐาน (Basic) ---
\examsection{ตอนที่ 1: พื้นฐาน — ระยะทาง, จุดกึ่งกลาง, ความชัน}

\problem{
  จงหาระยะทางระหว่างจุด $A(2, 3)$ และ $B(6, 6)$
}{
  $AB = \sqrt{(6-2)^2 + (6-3)^2} = \sqrt{16 + 9} = \sqrt{25} = 5$
}

\problem{
  จงหาจุดกึ่งกลางระหว่างจุด $P(-2, 5)$ และ $Q(4, -1)$
}{
  จุดกึ่งกลาง $= \left(\dfrac{-2+4}{2},\; \dfrac{5+(-1)}{2}\right) = (1,\; 2)$
}

\problem{
  จงหาความชันของเส้นตรงที่ผ่านจุด $(1, 2)$ และ $(4, 8)$
  พร้อมระบุว่าเส้นตรงนี้ชันขึ้นหรือชันลง
}{
  $m = \dfrac{8-2}{4-1} = \dfrac{6}{3} = 2$

  $m > 0$ แสดงว่าเส้นตรงชันขึ้น (increasing)
}

\problem{
  จงหาความชันของเส้นตรงที่ผ่านจุด $(3, 7)$ และ $(-2, 7)$
  และบอกชนิดของเส้นตรง
}{
  $m = \dfrac{7-7}{-2-3} = \dfrac{0}{-5} = 0$

  เส้นตรงขนานกับแกน $x$ (horizontal line)
}

% --- Part 2: การประยุกต์ (Intermediate) ---
\examsection{ตอนที่ 2: การประยุกต์ — สมการเส้นตรงและเวกเตอร์}

\problem{
  เส้นตรง $L_1$ มีความชัน $m_1 = \frac{2}{3}$
  และเส้นตรง $L_2$ มีความชัน $m_2 = -\frac{3}{2}$
  จงตรวจสอบว่า $L_1$ และ $L_2$ ขนานกันหรือตั้งฉากกัน
}{
  $m_1 \times m_2 = \dfrac{2}{3} \times \left(-\dfrac{3}{2}\right) = -1$

  ดังนั้น $L_1 \perp L_2$ (ตั้งฉากกัน)
}

\problem{
  จงหาสมการเส้นตรงที่ผ่านจุด $(3, 4)$ และมีความชัน $m = 2$
  ในรูป $y = mx + c$
}{
  $y - 4 = 2(x - 3)$

  $y - 4 = 2x - 6$

  $y = 2x - 2$
}

\problem{
  จงหาสมการเส้นตรงที่ผ่านจุด $A(-1, 3)$ และ $B(2, -6)$
  ในรูป $y = mx + c$
}{
  $m = \dfrac{-6-3}{2-(-1)} = \dfrac{-9}{3} = -3$

  $y - 3 = -3(x - (-1))$

  $y - 3 = -3x - 3$

  $y = -3x$
}

\problem{
  กำหนด $\vec{u} = (2, -3)$ และ $\vec{v} = (-5, 4)$
  จงหา $\vec{u} + \vec{v}$ และ $\vec{u} - \vec{v}$
}{
  $\vec{u} + \vec{v} = (2 + (-5),\; -3 + 4) = (-3,\; 1)$

  $\vec{u} - \vec{v} = (2 - (-5),\; -3 - 4) = (7,\; -7)$
}

\problem{
  จงหาขนาด (magnitude) ของเวกเตอร์ $\vec{v} = (-6, 8)$
}{
  $|\vec{v}| = \sqrt{(-6)^2 + 8^2} = \sqrt{36 + 64} = \sqrt{100} = 10$
}

% --- Part 3: ระดับ สอวน. (POSN Level) ---
\examsection{ตอนที่ 3: ระดับข้อสอบ สอวน.}

\problem{
  (ข้อสอบจริงปี 66 ข้อ 10)  เส้นตรงสองเส้นตั้งฉากกัน ณ จุดที่เส้นทั้งสองตัดแกน $X$ พอดี ถ้าสมการของเส้นตรงเส้นหนึ่งเป็น $3x-4y+5=0$ เส้นตรงอีกเส้นจะตัดแกน $Y$ ที่จุดใด 
}{
  $(0,-\dfrac{20}{9})$
}